\section{模型 OS 的产品结构：从聊天入口到任务入口}

把模型称为新一代 OS，最容易被误解成夸张说法。更准确地说，模型可能成为“任务入口”：用户不再先打开某个 app，而是先表达目标，再由模型选择工具、调用服务、组织上下文、执行动作并返回结果。这个转变会重塑应用分发、SaaS 形态和企业软件采购逻辑。

\subsection{三层产品结构}

\noindent\begin{longtable}{p{0.22\textwidth}p{0.34\textwidth}p{0.35\textwidth}}
\toprule
层级 & 职责 & 竞争焦点 \\
\midrule
模型层 & 理解意图、规划任务、生成代码或动作 & 基座能力、Coding 能力、工具使用能力。 \\
Harness 层 & 接入应用、权限、记忆、日志、重试、评估 & 工作流可靠性、成本、安全和可观测性。 \\
应用层 & 承载具体业务对象和用户关系 & 数据、场景、分发、用户习惯和付费。 \\
\bottomrule
\end{longtable}

\begin{knowledgebox}{为什么 OS 化会威胁应用入口}
如果用户通过模型表达目标，而模型负责选择应用和调用服务，传统 app 的入口价值会下降。应用仍然重要，但可能从“用户主动打开”变成“模型按任务调度”。
\end{knowledgebox}

\subsection{企业软件的变化}

企业软件过去围绕人类操作界面设计：表单、按钮、流程、权限和报表。Agent 时代，企业软件还需要提供可被模型调用的 API、可读状态、可审计日志和可恢复操作。谁能更早把产品变成 agent-friendly workflow，谁就更容易进入模型 OS 的调度网络。

\begin{importantbox}{Agent-friendly 软件的特征}
清晰 API、明确权限、可验证状态、可回滚动作、结构化日志、低成本沙盒和良好文档。没有这些，模型再强也难以安全稳定地操作企业系统。
\end{importantbox}

\subsection{本章小结}

模型 OS 不是替代操作系统内核，而是成为任务调度层。未来软件竞争会从“谁拥有界面入口”转向“谁能被模型可靠调用”。

